\chapter{Spectral Theory on Normed Spaces}

Throughout the spectral portions of this chapter, scalar fields are complex
unless explicitly stated otherwise.  For a linear operator
$T\colon\Dom(T)\subseteq X\to X$ and $\lambda\in\mathbb C$, write
\[
  T_\lambda:=T-\lambda I.
\]
When $T_\lambda$ is one-to-one, its inverse is understood as a map from
$\Ran(T_\lambda)$ to $\Dom(T)$.

\section{Resolvents and spectra}

\begin{notedefinition}[Regular values, resolvent, and spectrum]
Let $X\ne\{0\}$ be a complex Banach space and let
$T\colon\Dom(T)\subseteq X\to X$ be a closed linear operator.  A number
$\lambda\in\mathbb C$ is a \emph{regular value} if
$T_\lambda\colon\Dom(T)\to X$ is bijective and its inverse
\[
  R_\lambda(T):=T_\lambda^{-1}\colon X\to X
\]
is bounded.  The regular values form the \emph{resolvent set} $\rho(T)$,
and
\[
  \sigma(T):=\mathbb C\setminus\rho(T)
\]
is the \emph{spectrum}.

The spectrum is divided as follows:
\begin{itemize}
  \item $\lambda$ lies in the point spectrum $\sigma_p(T)$ when
        $T_\lambda$ is not injective; these are precisely the eigenvalues;
  \item it lies in the continuous spectrum when $T_\lambda$ is injective
        and has dense, non-surjective range;
  \item it lies in the residual spectrum when $T_\lambda$ is injective but
        its range is not dense.
\end{itemize}
The eigenspace corresponding to $\lambda$ is $\Ker(T_\lambda)$.
\end{notedefinition}

\begin{noteremark}[The convention used in the source]
The source calls $\lambda$ regular when $T_\lambda$ is injective and its
inverse is bounded on a dense range.  For a closed operator on a Banach space,
that convention is equivalent to the standard one above, as the next lemma
shows.  Bounded everywhere-defined operators are closed, so no distinction
arises in the remainder of these notes.
\end{noteremark}

\begin{notelemma}[A bounded inverse on a dense range]
Let $X$ be a Banach space, let
$T\colon\Dom(T)\subseteq X\to X$ be closed, and suppose $T_\lambda$ is
injective with dense range.  If
$T_\lambda^{-1}\colon\Ran(T_\lambda)\to X$ is bounded, then
$\Ran(T_\lambda)=X$.
\end{notelemma}

\begin{proof}
The inverse is closed because its graph is obtained from
$\Graph(T_\lambda)$ by interchanging the coordinates.  Its domain is also
closed: if $y_n\in\Ran(T_\lambda)$ and $y_n\to y$, boundedness makes
$T_\lambda^{-1}y_n$ Cauchy, hence convergent in $X$, and closedness of the
inverse puts $y$ back in its domain.  The range is therefore both dense and
closed, so it is all of $X$.
\end{proof}

\begin{noteexample}[A diagonal operator on $\ell^2$]
Let $(\alpha_j)$ be dense in $[0,1]$ and define
\[
  T(\xi_j)=(\alpha_j\xi_j)
  \qquad\text{on }\ell^2.
\]
Then
\[
  \sigma_p(T)=\{\alpha_j:j\ge1\},
  \qquad
  \sigma(T)=\overline{\{\alpha_j:j\ge1\}}=[0,1].
\]
If $\lambda\in[0,1]\setminus\{\alpha_j\}$, choose $j_n$ with
$\alpha_{j_n}\to\lambda$.  Then
\[
  R_\lambda(T)e_{j_n}
  =\frac1{\alpha_{j_n}-\lambda}e_{j_n},
\]
so the inverse is unbounded and $\lambda$ lies in the continuous spectrum.
\end{noteexample}

\begin{notetheorem}[Neumann series]
If $T\in\Bop{X}{X}$ on a Banach space and $\|T\|<1$, then
\[
  (I-T)^{-1}=\sum_{n=0}^{\infty}T^n
\]
with convergence in operator norm.
\end{notetheorem}

\begin{proof}
The series converges absolutely in the Banach algebra $\Bop{X}{X}$.  Its
partial sums satisfy
\[
  (I-T)\sum_{j=0}^{N}T^j=I-T^{N+1}
  =\left(\sum_{j=0}^{N}T^j\right)(I-T),
\]
and $T^{N+1}\to0$ in norm.
\end{proof}

\begin{notetheorem}[Openness of the resolvent and local expansion]
Let $T\in\Bop{X}{X}$ on a complex Banach space.  The resolvent set is open.
If $\lambda_0\in\rho(T)$ and
\[
  |\lambda-\lambda_0|<\frac1{\|R_{\lambda_0}(T)\|},
\]
then $\lambda\in\rho(T)$ and
\[
  R_\lambda(T)
  =\sum_{j=0}^{\infty}
     (\lambda-\lambda_0)^jR_{\lambda_0}(T)^{j+1}.
\]
Consequently $\sigma(T)$ is closed and the resolvent is operator-valued
holomorphic on $\rho(T)$.
\end{notetheorem}

\begin{proof}
Factor
\[
  T-\lambda I
  =(T-\lambda_0I)
   \bigl[I-(\lambda-\lambda_0)R_{\lambda_0}(T)\bigr]
\]
and apply the Neumann series to the second factor.
\end{proof}

\begin{notedefinition}[Spectral radius]
For $T\in\Bop{X}{X}$, the spectral radius is
\[
  r(T):=\sup\{ |\lambda|:\lambda\in\sigma(T)\}.
\]
\end{notedefinition}

\begin{notetheorem}[Resolvent identity and commutation]
For $\lambda,\mu\in\rho(T)$,
\[
  R_\mu(T)-R_\lambda(T)
  =(\mu-\lambda)R_\mu(T)R_\lambda(T).
\]
If $S\in\Bop{X}{X}$ commutes with $T$, then it commutes with every
$R_\lambda(T)$.  In particular, all resolvents commute with one another.
\end{notetheorem}

\begin{proof}
The first formula follows from
\[
  A^{-1}-B^{-1}=A^{-1}(B-A)B^{-1}
\]
with $A=T-\mu I$ and $B=T-\lambda I$.  If $ST=TS$, then
$S(T-\lambda I)=(T-\lambda I)S$; multiplying by the inverse on both sides
gives $SR_\lambda=R_\lambda S$.
\end{proof}

\begin{notetheorem}[Polynomial spectral mapping]
For a polynomial $p$ and $T\in\Bop{X}{X}$,
\[
  \sigma(p(T))=p(\sigma(T)).
\]
\end{notetheorem}

\begin{proof}
If $\mu\notin p(\sigma(T))$, factor
\[
  p(z)-\mu=c\prod_{j=1}^{m}(z-\lambda_j).
\]
Every $\lambda_j$ belongs to $\rho(T)$, so the commuting product
$p(T)-\mu I$ is invertible.

Conversely, if $\mu=p(\lambda)$ with $\lambda\in\sigma(T)$, polynomial
division gives
\[
  p(T)-p(\lambda)I=(T-\lambda I)q(T).
\]
If the left-hand side were invertible, the commuting factors would give both
a left and a right inverse for $T-\lambda I$, contradicting
$\lambda\in\sigma(T)$.
\end{proof}

\begin{notelemma}[Eigenvectors for distinct eigenvalues]
Eigenvectors corresponding to distinct eigenvalues are linearly independent.
\end{notelemma}

\begin{noteexample}[Idempotent operators]
If $T^2=T$ and $T\ne0,I$, then
\[
  \sigma(T)=\{0,1\}.
\]
Indeed, polynomial spectral mapping applied to $p(z)=z^2-z$ gives
$\sigma(T)\subseteq\{0,1\}$.  Since $T$ and $T-I$ are both noninjective,
$0$ and $1$ are eigenvalues.
\end{noteexample}

\begin{notetheorem}[Resolvent norm and distance to the spectrum]
For $\lambda\in\rho(T)$,
\[
  \|R_\lambda(T)\|
  \ge \frac1{\dist(\lambda,\sigma(T))}.
\]
Hence the resolvent norm tends to infinity as $\lambda$ approaches the
spectrum through $\rho(T)$.
\end{notetheorem}

\begin{proof}
The local expansion shows that the open disk of radius
$1/\|R_\lambda(T)\|$ centred at $\lambda$ lies in $\rho(T)$.  Its radius
cannot exceed the distance from $\lambda$ to the complementary set
$\sigma(T)$.
\end{proof}

\begin{notetheorem}[Nonempty compact spectrum]
If $X\ne\{0\}$ is a complex Banach space and $T\in\Bop{X}{X}$, then
$\sigma(T)$ is nonempty and compact, and
\[
  \sigma(T)\subseteq\{\lambda:|\lambda|\le\|T\|\}.
\]
\end{notetheorem}

\begin{proof}
For $|\lambda|>\|T\|$ the Neumann expansion
\[
  R_\lambda(T)
  =-\frac1\lambda\sum_{n=0}^{\infty}\frac{T^n}{\lambda^n}
\]
exists, so the spectrum is bounded.  It is closed by openness of the
resolvent.

If the spectrum were empty, then for every $x\in X$ and $f\in X^*$ the
scalar function $h(\lambda)=f(R_\lambda(T)x)$ would be entire.  The displayed
expansion shows that $h(\lambda)\to0$ as $|\lambda|\to\infty$, so Liouville's
theorem makes $h$ identically zero.  Since bounded functionals separate
points, $R_\lambda(T)=0$, impossible for an inverse.
\end{proof}

\begin{notetheorem}[Spectral radius formula]
For $T\in\Bop{X}{X}$ on a complex Banach space,
\[
  r(T)=\lim_{n\to\infty}\|T^n\|^{1/n}
      =\inf_{n\ge1}\|T^n\|^{1/n}.
\]
\end{notetheorem}

\begin{proof}
The sequence $a_n=\|T^n\|$ is submultiplicative, so
$a_n^{1/n}$ has a limit equal to its infimum.  Polynomial spectral mapping
gives
\[
  r(T)^n=r(T^n)\le\|T^n\|,
\]
therefore $r(T)\le\lim_n\|T^n\|^{1/n}$.

For the reverse inequality, fix $\rho>r(T)$.  The circle
$|\lambda|=\rho$ lies in the resolvent set.  The Laurent coefficients of the
resolvent at infinity satisfy the Cauchy formula
\[
  T^n=-\frac1{2\pi i}\int_{|\lambda|=\rho}
       \lambda^nR_\lambda(T)\,\mathrm d\lambda.
\]
Thus, with $M_\rho=\max_{|\lambda|=\rho}\|R_\lambda(T)\|$,
\[
  \|T^n\|\le M_\rho\rho^{n+1}.
\]
Taking $n$th roots and then letting $\rho\downarrow r(T)$ gives the reverse
inequality.
\end{proof}

\section{Compact operators}

\begin{notedefinition}[Compact linear operator]
Let $X,Y$ be normed spaces.  A linear operator $T\colon X\to Y$ is
\emph{compact} if it maps every bounded subset of $X$ to a relatively compact
subset of $Y$.
\end{notedefinition}

\begin{notetheorem}[Basic compactness criteria]
Let $T\colon X\to Y$ be linear.
\begin{enumerate}
  \item $T$ is compact if and only if every bounded sequence $(x_n)$ has a
        subsequence $(x_{n_k})$ for which $(Tx_{n_k})$ converges in $Y$.
  \item Every compact linear operator is bounded.
  \item The identity on an infinite-dimensional normed space is not compact.
  \item If $T$ is bounded and either its domain or its range is
        finite-dimensional, then $T$ is compact.
\end{enumerate}
\end{notetheorem}

\begin{proof}
The first assertion is sequential compactness in metric spaces.  For the
second, the relatively compact set $T(B_X)$ is bounded, so
$\sup_{\|x\|\le1}\|Tx\|<\infty$.  The third follows because the closed unit
ball is not compact in infinite dimension.  The final assertion follows from
finite-dimensional Heine--Borel compactness.
\end{proof}

\begin{notetheorem}[Norm limits of compact operators]
Let $Y$ be Banach.  If $T_n\in\Kop{X}{Y}$ and $\|T_n-T\|\to0$, then $T$ is
compact.  Consequently $\Kop{X}{Y}$ is a closed subspace of
$\Bop{X}{Y}$.
\end{notetheorem}

\begin{proof}
Given a bounded sequence, successively extract subsequences on which
$T_1,T_2,\ldots$ converge, and take the diagonal subsequence $(x_{n_k})$.
Put $M=\max\{1,\sup_k\|x_{n_k}\|\}$.  Choose $j$ with
$\|T-T_j\|<\varepsilon/(3M)$ and then $k,\ell$ large enough that
$\|T_jx_{n_k}-T_jx_{n_\ell}\|<\varepsilon/3$.  It follows that
$(Tx_{n_k})$ is Cauchy, hence convergent in $Y$.
\end{proof}

\begin{notetheorem}[Weak-to-strong continuity]
If $T\colon X\to Y$ is compact and $x_n\weakto x$, then
\[
  Tx_n\longrightarrow Tx
\]
in norm.
\end{notetheorem}

\begin{proof}
The sequence $(x_n)$ is bounded.  Every subsequence of $(Tx_n)$ has a
norm-convergent further subsequence by compactness.  Its norm limit must be
$Tx$, because bounded linear maps preserve weak convergence and weak limits
are unique.  Therefore the whole sequence converges to $Tx$ in norm.
\end{proof}

\begin{notedefinition}[$\varepsilon$-nets and total boundedness]
A finite set $F$ is an $\varepsilon$-net for $B$ if every point of $B$ lies
within distance $\varepsilon$ of a point of $F$.  A set is \emph{totally
bounded} if it has a finite $\varepsilon$-net for every $\varepsilon>0$.
\end{notedefinition}

\begin{notelemma}[Total boundedness]
For a subset $B$ of a metric space:
\begin{enumerate}
  \item relative compactness implies total boundedness;
  \item total boundedness in a complete space implies relative compactness;
  \item the centres of a finite $\varepsilon$-net may be chosen in $B$;
  \item every totally bounded set is separable.
\end{enumerate}
\end{notelemma}

\begin{proof}
A compact set is covered by finitely many $\varepsilon$-balls.  Conversely,
successive finite $2^{-n}$-nets yield a Cauchy subsequence of every sequence;
completeness supplies its limit.  Replacing each nonempty covering ball by a
point of $B$ proves the third assertion.  The union of finite $1/n$-nets is a
countable dense subset.
\end{proof}

\begin{notetheorem}[Separable range and extension to the completion]
Let $T\colon X\to Y$ be compact.
\begin{enumerate}
  \item $\Ran(T)$ is separable.
  \item If $Y$ is Banach and $\widehat X$ is the completion of $X$, the
        bounded extension $\widehat T\colon\widehat X\to Y$ is compact.
\end{enumerate}
\end{notetheorem}

\begin{proof}
Since $X=\bigcup_{n\ge1}nB_X$ and each $T(nB_X)$ is relatively compact and
therefore separable, a countable union of countable dense sets is dense in
$\Ran(T)$.

For the extension, let $(\widehat x_n)$ be bounded in $\widehat X$ and choose
$x_n\in X$ with $\|x_n-\widehat x_n\|\to0$.  The sequence $(x_n)$ is bounded,
so some $Tx_{n_k}$ converges.  Since
$\widehat T(\widehat x_{n_k}-x_{n_k})\to0$, the corresponding
$\widehat T\widehat x_{n_k}$ also converges.
\end{proof}

\begin{notetheorem}[Schauder theorem]
If $T\colon X\to Y$ is compact, then its adjoint
$T^*\colon Y^*\to X^*$ is compact.
\end{notetheorem}

\begin{proof}
The set $K=\overline{T(B_X)}$ is compact and hence totally bounded.  Given
$\varepsilon>0$, choose $y_1,\ldots,y_N\in K$ so that every $Tx$ with
$\|x\|\le1$ is within $\varepsilon/4$ of some $y_j$.  On the dual unit ball,
the coordinate map
\[
  g\longmapsto(g(y_1),\ldots,g(y_N))
\]
has totally bounded image in the finite-dimensional space $\mathbb K^N$.
Choose finitely many representatives so that any two functionals in the same
coordinate cell differ by less than $\varepsilon/2$ on each $y_j$.
Then, for $\|x\|\le1$ and the corresponding $j$,
\[
 |(g-h)(Tx)|
 \le |(g-h)(y_j)|+\|g-h\|\,\|Tx-y_j\|<\varepsilon.
\]
Thus $T^*(B_{Y^*})$ is totally bounded in the Banach space $X^*$, hence
relatively compact.
\end{proof}

\section{Spectral structure of compact operators}

\begin{notetheorem}[Nonzero eigenvalues]
Let $T\colon X\to X$ be compact.  Its set of eigenvalues is at most
countable, and $0$ is its only possible accumulation point.  Equivalently,
for every $c>0$ there are only finitely many eigenvalues satisfying
$|\lambda|\ge c$.
\end{notetheorem}

\begin{proof}
Suppose distinct eigenvalues $\lambda_n$ satisfy $|\lambda_n|\ge c>0$, and
choose eigenvectors so that
$M_n=\Span\{x_1,\ldots,x_n\}$ is strictly increasing.  By Riesz's lemma choose
$y_n\in M_n$ with $\|y_n\|=1$ and
$\dist(y_n,M_{n-1})\ge1/2$.  Each $M_n$ is $T$-invariant.  For $m<n$,
\[
  Ty_n-Ty_m
  =\lambda_n y_n+z_{m,n}
  \qquad(z_{m,n}\in M_{n-1}),
\]
so $\|Ty_n-Ty_m\|\ge c/2$.  This contradicts compactness.  The finiteness
claim implies countability by taking $c=1/k$.
\end{proof}

\begin{notelemma}[Composition]
If $T$ is compact and $S$ is bounded, then $TS$ and $ST$ are compact whenever
the compositions are defined.
\end{notelemma}

\begin{notetheorem}[Finite-dimensional null spaces]
Let $T\colon X\to X$ be compact and $\lambda\ne0$.  Then
\[
  \dim\Ker(T-\lambda I)<\infty.
\]
More generally, every $\Ker(T-\lambda I)^n$ is finite-dimensional.
\end{notetheorem}

\begin{proof}
On $N=\Ker(T-\lambda I)$ one has $Tx=\lambda x$.  Hence the image under $T$
of the closed unit ball of $N$ is a scalar multiple of that ball.  Compactness
of $T$ makes the ball compact, so $N$ is finite-dimensional.

For powers, the binomial expansion gives
\[
  (T-\lambda I)^n=(-\lambda)^nI+K_n,
\]
where $K_n$ is compact.  The same argument applied to $K_n$ shows that the
kernel is finite-dimensional.
\end{proof}

\begin{notetheorem}[Closed range of a compact perturbation of a scalar]
Let $X$ be Banach, let $T\colon X\to X$ be compact, and let
$\lambda\ne0$.  Then
\[
  \Ran(T-\lambda I)
\]
is closed.  The same is true for every power $(T-\lambda I)^n$.
\end{notetheorem}

\begin{proof}
Put $A=T-\lambda I$ and $N=\Ker A$.  The induced injective map
\[
  \widetilde A\colon X/N\to X,
  \qquad \widetilde A(x+N)=Ax,
\]
is bounded below.  Otherwise there are cosets of norm one represented by a
bounded sequence $(x_n)$ with $Ax_n\to0$.  Compactness gives a subsequence
for which $Tx_n$ converges, and then
\[
  \lambda x_n=Tx_n-Ax_n
\]
converges.  Its limit belongs to $N$, contradicting
$\dist(x_n,N)=1$.  A bounded-below operator has closed range.

For $A^n=(-\lambda)^nI+K_n$, apply the same argument to the compact operator
$K_n$.
\end{proof}

\begin{notecorollary}
Let $T\colon X\to Y$ be compact between Banach spaces.  Then
$\Ran(T)$ contains no closed infinite-dimensional subspace of $Y$.
\end{notecorollary}

\begin{proof}
If a closed infinite-dimensional subspace $M\subseteq\Ran(T)$ existed, the
map $T\colon T^{-1}(M)\to M$ would be a bounded surjection between Banach
spaces.  The open mapping theorem would make a fixed multiple of the image of
the unit ball contain $B_M$.  That would make $B_M$ relatively compact,
contradicting infinite dimensionality.
\end{proof}

\section{Stabilization and the Riesz--Schauder theorem}

Throughout this section, $X$ is a Banach space and
$A=T-\lambda I$, where $T\colon X\to X$ is compact and $\lambda\ne0$.
Write
\[
  N_n:=\Ker A^n,
  \qquad
  R_n:=\Ran A^n.
\]
Then $N_0\subseteq N_1\subseteq\cdots$ and
$X=R_0\supseteq R_1\supseteq\cdots$.

\begin{notelemma}[Stabilization of null spaces]
There is a least integer $p\ge0$ such that
\[
  N_p=N_{p+1}=N_{p+2}=\cdots.
\]
If $p>0$, all inclusions $N_0\subsetneq N_1\subsetneq\cdots\subsetneq N_p$
are proper.
\end{notelemma}

\begin{proof}
If the chain were strictly increasing forever, Riesz's lemma would provide
$x_n\in N_n$ with $\|x_n\|=1$ and
$\dist(x_n,N_{n-1})\ge1/2$.  For $m<n$,
$Ax_n\in N_{n-1}$ and $Tx_m=Ax_m+\lambda x_m\in N_m\subseteq N_{n-1}$.
Thus
\[
  Tx_n-Tx_m=\lambda x_n+z_{m,n}
  \qquad(z_{m,n}\in N_{n-1}),
\]
so $\|Tx_n-Tx_m\|\ge|\lambda|/2$, contradicting compactness.  Once
$N_p=N_{p+1}$, applying $A$ shows inductively that every later equality also
holds.
\end{proof}

\begin{notelemma}[Stabilization of ranges]
There is a least integer $q\ge0$ such that
\[
  R_q=R_{q+1}=R_{q+2}=\cdots.
\]
If $q>0$, all inclusions $R_0\supsetneq R_1\supsetneq\cdots\supsetneq R_q$
are proper.
\end{notelemma}

\begin{proof}
The adjoint $T^*$ is compact and
\[
  A^*=T^*-\lambda I.
\]
The null-space chain of $A^*$ stabilizes.  Since every $R_n$ is closed and
\[
  R_n^\perp=\Ker(A^*)^n,
\]
equality of two successive adjoint kernels is equivalent, by Hahn--Banach,
to equality of the corresponding ranges.
\end{proof}

\begin{notetheorem}[Equal ascent and descent]
The least stabilization indices above are equal: $p=q=:r$.  Consequently
\[
  \Ker A^r=\Ker A^{r+1}=\cdots,
  \qquad
  \Ran A^r=\Ran A^{r+1}=\cdots.
\]
If $r>0$, every preceding inclusion in both chains is proper.
\end{notetheorem}

\begin{proof}
If the null chain stabilizes at $p$, then
\[
  \Ker A^p\cap\Ran A^p=\{0\}.
\]
Indeed, if $x=A^py$ and $A^px=0$, then $A^{2p}y=0$; stabilization gives
$y\in\Ker A^p$, hence $x=0$.

If the range chain stabilizes at $q$, then
\[
  X=\Ker A^q+\Ran A^q.
\]
For $x\in X$, equality $\Ran A^q=\Ran A^{2q}$ gives
$A^qx=A^{2q}y$ for some $y$, and then
$x-A^qy\in\Ker A^q$.

If $p\le q$, then $\Ker A^p=\Ker A^q$.  Applying $A^p$ to
$X=\Ker A^q+\Ran A^q$ gives
$\Ran A^p=\Ran A^{p+q}=\Ran A^q$, so $q\le p$.  If $q\le p$, decompose every
$x\in\Ker A^p$ as $x=u+v$ with $u\in\Ker A^q$ and
$v\in\Ran A^q=\Ran A^p$.  Then
$v\in\Ker A^p\cap\Ran A^p=\{0\}$, so
$\Ker A^p=\Ker A^q$ and $p\le q$.  Hence $p=q$.
\end{proof}

\begin{notetheorem}[Nonzero spectral values are eigenvalues]
Let $X$ be a complex Banach space and $T\in\Kop{X}{X}$.  Every
$\lambda\in\sigma(T)\setminus\{0\}$ is an eigenvalue.
\end{notetheorem}

\begin{proof}
If $A=T-\lambda I$ were injective, its null-space chain would stabilize at
$r=0$.  Equal ascent and descent would then give
$\Ran A=X$.  Thus $A$ would be a bounded bijection of Banach spaces, and the
bounded inverse theorem would put $\lambda$ in the resolvent, a contradiction.
\end{proof}

\begin{notetheorem}[Riesz decomposition]
With $r$ as above,
\[
  X=\Ker A^r\oplus\Ran A^r.
\]
Both summands are closed and invariant under $A$ and $T$.
\end{notetheorem}

\begin{proof}
The intersection is zero by the argument in the preceding theorem, and the
sum is all of $X$ by range stabilization.  Closedness of the two summands was
proved above.
\end{proof}

\begin{noteexercise}[Two noncompact operators]
\begin{enumerate}
  \item For the diagonal operator on $\ell^2$ with $(\alpha_j)$ dense in
        $[0,1]$, choose a subsequence with $|\alpha_{j_k}|\ge1/2$.  Then
        $\|Te_{j_k}-Te_{j_\ell}\|\ge1/\sqrt2$, so $T$ is not compact.
  \item On $C[0,1]$, the multiplication operator $(Tx)(t)=tx(t)$ is not
        compact.  The bounded sequence $x_n(t)=t^n$ has images
        $Tx_n=t^{n+1}$.  Every subsequence converges pointwise to the
        discontinuous function that is $0$ on $[0,1)$ and $1$ at $1$, so no
        subsequence can converge uniformly.
\end{enumerate}
\end{noteexercise}

\section{Fredholm solvability}

Throughout this section, $X$ is a complex Banach space.  Let
$T\in\Kop{X}{X}$, let $\lambda\ne0$, and put
\[
  A:=T-\lambda I.
\]
For both real and complex Banach spaces,
\[
  A^*=T^*-\lambda I.
\]
There is no conjugation here: the Banach-space adjoint acts on complex-linear
functionals, so $(\lambda I)^*=\lambda I$.  Conjugation appears only after
identifying a Hilbert space with its dual through the Riesz map.

\begin{notetheorem}[Solvability of $Ax=y$]
The equation
\[
  Ax=y
\]
has a solution if and only if
\[
  f(y)=0
  \qquad\text{for every }f\in\Ker A^*.
\]
In other words,
\[
  \Ran A={}^{\perp}(\Ker A^*).
\]
\end{notetheorem}

\begin{proof}
The identity $(\Ran A)^\perp=\Ker A^*$ holds for every bounded operator.  Here
$\Ran A$ is closed, so the Hahn--Banach annihilator theorem gives
\[
  \Ran A={}^{\perp}\bigl((\Ran A)^\perp\bigr)
        ={}^{\perp}(\Ker A^*).
\]
\end{proof}

\begin{notedefinition}[Normally solvable equation]
A bounded operator equation $Bx=y$ is called \emph{normally solvable} when
\[
  y\in\Ran B
  \quad\Longleftrightarrow\quad
  f(y)=0\text{ for every }f\in\Ker B^*.
\]
Equivalently, $\Ran B$ is closed.
\end{notedefinition}

\begin{notelemma}[A bounded choice of solution]
There is $C>0$ such that whenever $y\in\Ran A$, one can choose a solution
$\widetilde x$ of $A\widetilde x=y$ satisfying
\[
  \|\widetilde x\|\le C\|y\|.
\]
\end{notelemma}

\begin{proof}
The finite-dimensional space $\Ker A$ has a closed complement $M$.  The
restriction
\[
  A|_M\colon M\to\Ran A
\]
is a bounded bijection between Banach spaces.  Its inverse is bounded, and
$\widetilde x=(A|_M)^{-1}y$ has the required estimate.
\end{proof}

\begin{notetheorem}[Solvability of the adjoint equation]
For $g\in X^*$, the equation
\[
  A^*f=g
\]
has a solution if and only if
\[
  g(x)=0
  \qquad\text{for every }x\in\Ker A.
\]
Equivalently,
\[
  \Ran A^*=(\Ker A)^\perp.
\]
\end{notetheorem}

\begin{proof}
Define on $\Ran A$ the functional $F(Ax)=g(x)$.  The condition on $g$
makes $F$ well defined.  For $y=Ax$, choose a solution $\widetilde x$ with
$\|\widetilde x\|\le C\|y\|$; then
\[
  |F(y)|=|g(\widetilde x)|\le C\|g\|\|y\|.
\]
Thus $F$ is bounded.  Extend it to $f\in X^*$ by Hahn--Banach.  For every
$x\in X$,
$(A^*f)(x)=f(Ax)=F(Ax)=g(x)$, so $A^*f=g$.
\end{proof}

\begin{notetheorem}[Injectivity, surjectivity, and bounded inverses]
The following are equivalent:
\begin{enumerate}
  \item $A$ is injective;
  \item $A$ is surjective;
  \item $A$ is bijective.
\end{enumerate}
When these conditions hold, $A^{-1}$ is bounded.  The same equivalence holds
for $A^*$.
\end{notetheorem}

\begin{proof}
If $A$ is injective, its ascent is zero; equal ascent and descent imply that
its range chain stabilizes at zero, so $\Ran A=X$.  Conversely, if $A$ is
surjective, then $R_0=R_1=X$, so the descent is zero.  Equal ascent and descent
then gives $N_0=N_1$, hence $\Ker A=\{0\}$.  A bounded bijection between
Banach spaces has bounded inverse.  Applying the same reasoning to the compact
operator $T^*$ gives the adjoint statement.
\end{proof}

\begin{notelemma}[Biorthogonal system]
If $f_1,\ldots,f_m\in X^*$ are linearly independent, there are
$z_1,\ldots,z_m\in X$ such that
\[
  f_j(z_k)=\delta_{jk}.
\]
\end{notelemma}

\begin{proof}
Proceed by induction.  Suppose $z_1,\ldots,z_m$ have been constructed.  The
functional $f_{m+1}$ cannot vanish on
$\bigcap_{j=1}^{m}\Ker f_j$, for otherwise it would be a linear combination
of $f_1,\ldots,f_m$.  Choose $w$ in that intersection with
$f_{m+1}(w)\ne0$ and put $z_{m+1}=w/f_{m+1}(w)$.  Replacing each old $z_j$ by
$z_j-f_{m+1}(z_j)z_{m+1}$ preserves the previous relations and makes the new
one biorthogonal as well.
\end{proof}

\begin{notetheorem}[Equality of nullities]
For $A=T-\lambda I$ with $\lambda\ne0$,
\[
  \dim\Ker A=\dim\Ker A^*<\infty.
\]
Equivalently, $A$ has Fredholm index zero.
\end{notetheorem}

\begin{proof}
Let $r$ be the common ascent and descent.  The Riesz decomposition gives
\[
  X=N\oplus R,
  \qquad N=\Ker A^r,
  \quad R=\Ran A^r.
\]
The restriction of $A$ to $R$ is bijective, while $N$ is finite-dimensional
and $A|_N$ is nilpotent.  Therefore finite-dimensional rank--nullity gives
\[
  \operatorname{codim}\Ran A=\dim\Ker A.
\]
Because $\Ran A$ is closed,
$\Ker A^*=(\Ran A)^\perp$, whose dimension equals
$\dim(X/\Ran A)$.  This proves the equality.
\end{proof}

\begin{notecorollary}
For a compact operator on a complex Banach space and $\lambda\ne0$,
\[
  \lambda\in\sigma_p(T)
  \quad\Longleftrightarrow\quad
  \lambda\in\sigma_p(T^*).
\]
In particular, every nonzero spectral value of $T$ is an eigenvalue.
\end{notecorollary}

\section{The Fredholm alternative and integral equations}

\begin{notedefinition}[Fredholm alternative]
A bounded operator $B\colon X\to X$ is said to satisfy the Fredholm
alternative if exactly one of the following cases occurs:
\begin{enumerate}
  \item $B$ and $B^*$ are bijective, so $Bx=y$ and $B^*f=g$ have
        unique solutions for every $y\in X$ and $g\in X^*$;
  \item the two homogeneous equations have finite-dimensional nonzero
        solution spaces of the same dimension, and the nonhomogeneous
        equations are solvable precisely under the orthogonality conditions
        \[
          f(y)=0\quad(f\in\Ker B^*),
          \qquad
          g(x)=0\quad(x\in\Ker B).
        \]
\end{enumerate}
\end{notedefinition}

\begin{notetheorem}[Fredholm alternative for compact operators]
If $T\colon X\to X$ is compact and $\lambda\ne0$, then
$T-\lambda I$ satisfies the Fredholm alternative.
\end{notetheorem}

\begin{proof}
This is the combination of closed range, the two solvability criteria,
injectivity--surjectivity equivalence, and equality of nullities established
above.
\end{proof}

\begin{notetheorem}[Arzel\`a--Ascoli, sequential form]
A uniformly bounded and equicontinuous sequence in $C([a,b])$ has a uniformly
convergent subsequence.
\end{notetheorem}

\begin{notetheorem}[Compact integral operator]
Let $k\in C([a,b]\times[a,b])$ and define
\[
  (Tx)(s)=\int_a^b k(s,t)x(t)\,\mathrm dt
  \qquad(x\in C([a,b])).
\]
Then $T\in\Kop{C([a,b])}{C([a,b])}$.
\end{notetheorem}

\begin{proof}
Linearity is immediate, and
\[
  \|Tx\|_\infty
  \le \|x\|_\infty
      \max_{s\in[a,b]}\int_a^b|k(s,t)|\,\mathrm dt.
\]
Let $(x_n)$ be bounded by $M$ and put $y_n=Tx_n$.  Uniform continuity of $k$
on the compact square implies that, whenever $|s_1-s_2|$ is sufficiently
small,
\[
  |y_n(s_1)-y_n(s_2)|
  \le M\int_a^b|k(s_1,t)-k(s_2,t)|\,\mathrm dt
  <\varepsilon
\]
uniformly in $n$.  Thus $(y_n)$ is uniformly bounded and equicontinuous, so
Arzel\`a--Ascoli provides a uniformly convergent subsequence.
\end{proof}

\begin{notecorollary}[Fredholm integral equation]
For $\mu\ne0$, consider
\[
  x(s)-\mu\int_a^b k(s,t)x(t)\,\mathrm dt=\widetilde y(s).
\]
Exactly one of the following occurs:
\begin{enumerate}
  \item for every $\widetilde y\in C([a,b])$ there is a unique solution;
  \item the homogeneous equation has a finite-dimensional nonzero solution
        space, and the nonhomogeneous equation is solvable exactly when its
        right-hand side annihilates every solution of the adjoint homogeneous
        equation in the dual space.
\end{enumerate}
On $L^2([a,b])$, when $k\in L^2([a,b]^2)$, the same conclusion follows from
compactness of the Hilbert--Schmidt operator; its Hilbert-space adjoint has
kernel $\overline{k(t,s)}$.
\end{notecorollary}
